\section{Introduction}

For \(n\ge0\), let
\(QS^n=\operatorname*{colim}_r\Omega^rS^{n+r}\), and let \(Q_0S^0\)
be its basepoint component. For \(d>0\), the identification
\(\pi_d(Q_0S^0)=\pi_d^S\) gives the mod-2 Hurewicz homomorphism
\[
h_d:\pi_d^S\longrightarrow H_d(Q_0S^0;\F_2).
\]
It sends a representative \(S^d\to Q_0S^0\) to the image of the
mod-2 fundamental class. Curtis's conjecture describes the positive-degree
image of these maps \cite{Curtis1975}.
The target has an explicit homology algebra, while the stable homotopy
groups in the source contain many families whose images are difficult
to determine. The image problem asks which of the readily described
homology classes are represented by maps from spheres.

Curtis proposed the Hopf--Kervaire description in his study of the
Dyer--Lashof and lambda algebras \cite{Curtis1975}. May's subsequent
account records a gap found by Wellington in the proposed proof
\cite[p.~484, footnote~3]{May1977}. We prove the following form of
the conjecture.

\begin{theorem}\label{thm:curtis}
The image of \(\bigoplus_{d>0}h_d\) is spanned over \(\F_2\) by
the Hurewicz images of the Hopf-invariant-one classes
\(\eta,\nu,\sigma\), in degrees \(1,3,7\), and the
Kervaire-invariant-one classes \(\theta_j\), in degrees
\(2^{j+1}-2\), \(j\ge1\), whenever they exist.
In particular, Curtis's conjecture is true in all positive degrees.
\end{theorem}

Here Kervaire invariant one means that the corresponding framed manifold,
under the Pontryagin--Thom identification, has Arf invariant one. The
notation \(\theta_j\) is used only when such a class exists. The theorem
does not posit existence in any additional degree.

Eccles's conjecture concerns the analogous realization problem for
\(QX\). In the formulation recorded in
\cite[pp.~14--15]{Zare2009Thesis}, for a path-connected CW complex
\(X\) of finite type, a positive-dimensional sphere
map to \(QX\) with nonzero mod-2 Hurewicz image has stable adjoint
detected either by homology or by a primary cohomology operation in
its mapping cone. Theorem~\ref{thm:curtis} gives this conclusion
when \(X\) is a positive-dimensional sphere.

For \(n>0\), write \(e_n\in H_n(QS^n;\F_2)\) for the bottom spherical
class, and write \(\sigma_*^n\) for \(n\)-fold homology suspension from
\(Q_0S^0\).

\begin{corollary}[Eccles's conjecture for spheres]\label{cor:eccles-spheres}
For every integer \(n>0\), each nonzero mod-2 spherical class in
\(H_*(QS^n;\F_2)\) is either \(e_n\), or a nonzero class of the form
\[
 \sigma_*^n h_d(\alpha),\qquad
 (\alpha,d)\in\{(\eta,1),(\nu,3),(\sigma,7)\}.
\]
Consequently Eccles's conjecture holds for every sphere \(S^n\), \(n>0\):
the alternatives are the bottom stably spherical class and
Hopf-invariant-one classes.
\end{corollary}
\begin{proof}
For a spherical representative \(f:S^{n+d}\to QS^n\) with \(d>0\),
adjunction gives \(f_0:S^d\to Q_0S^0\), and naturality gives
\(h(f)=\sigma_*^n h_d([f_0])\). By Theorem~\ref{thm:curtis}, the
level-zero image is spanned by the Hopf and Kervaire images. The latter
are squares killed by the first homology suspension, leaving the stated
Hopf possibilities; this is the Curtis-to-Eccles sphere implication
recorded in \cite[Chapter~5, pp.~81--82]{Zare2009Thesis}. In degree
\(n\), \(\pi_n(QS^n)=\pi_0^S=\Z\) gives the bottom image represented by
the unit \(S^n\to QS^n\); connectivity excludes lower degrees.
\end{proof}

\begin{corollary}[Finite positive stable Hurewicz image]\label{cor:finite-image}
The image of \(\bigoplus_{d>0}h_d\) is finite, and
\[
 h_d=0\qquad\text{if }d>0\text{ and }
 d\notin\{1,2,3,6,7,14,30,62,126\}.
\]
In particular, the positive stable mod-2 Hurewicz image vanishes above
degree \(126\).
\end{corollary}
\begin{proof}
Hill, Hopkins, and Ravenel prove that \(\theta_j\) can exist only for
\(j\le6\) \cite[Theorem~1.1]{HillHopkinsRavenel2016}. By Theorem~\ref{thm:curtis},
the positive image is therefore spanned by the named images in Hopf degrees
\(1,3,7\) and Kervaire degrees \(2,6,14,30,62,126\). Finitely many
generators over \(\F_2\) give a finite image and the stated support.
\end{proof}

Corollary~\ref{cor:eccles-spheres} concerns only positive-dimensional
spheres, not general spaces. The support in Corollary~\ref{cor:finite-image}
is an upper bound; it does not assert nonvanishing in every listed degree
or an exact dimension of the image.

Adams's Hopf invariant one theorem identifies the Hopf cases
\cite{Adams1960}. Browder's work relates the framed Kervaire invariant
to the classes \(h_j^2\) in the mod-2 Adams spectral sequence
\cite{Browder1969}. Hill--Hopkins--Ravenel's nonexistence theorem
leaves dimension \(126\) as the last possible Kervaire dimension.
In their preprint, Lin, Wang, and Xu prove existence in that dimension
\cite[Theorems~1.1 and~1.4]{LinWangXu2025}. This endpoint result is
separate from the upper bound used in Corollary~\ref{cor:finite-image}.

On the homology side, the free infinite-loop-space
calculation describes a large algebra generated by Dyer--Lashof
operations \cite{CLM1976}. A spherical class must be primitive and
annihilated by positive Steenrod operations, but these necessary
conditions alone do not establish its realization. The theorem concerns
the image of actual stable homotopy classes.

Eccles and Zare give restrictions on spherical squares and their roots,
including the exclusion of higher iterated squares
\cite[Theorem~3.10 and Lemma~3.12]{EcclesZare2009}.
Section~\ref{sec:powers} proves the exclusion of fourth powers by
spectrum pairings and integral characteristic classes; the conclusion
itself is already present in that prior work. The Dickson-algebra
approach to spherical classes also appears in Hung and Peterson's
study of the Lannes--Zarati homomorphism and products in the unstable
Adams \(E_2\)-term \cite{HungPeterson1998}.
Hunter's cohomological construction of the Curtis--Wellington spectral
sequence exhibits Dickson-algebra quotients in a width filtration
\cite[arXiv version~2, Theorem~1.2]{Hunter2025}.

The argument here combines classical detection by mapping-cone
operations with an obstruction for the specific Steenrod modules
dual to primitives. The cone of an \emph{actual} sphere map supplies
a nonzero extension. A detecting functional inherited from the
preceding suspension level then restricts that extension through the
algebraic obstruction proved in Section~\ref{sec:algebra}.

The free-algebra model constructed below is compatible with Kuhn's
structured splitting and completion maps \cite{Kuhn2006}; we prove
the required finite weight projections and their suspension
compatibility directly. For the final passage from leading
weight to Adams filtration, we use Kuhn's lifting theorem
\cite{Kuhn2018} and the derived ideal-power constructions of
\cite{KuhnPereira2017}.

\paragraph{Outline of the proof.}
The proof is organized around the last nonzero homology suspension of a
spherical class. For a representative \(f_0:S^d\to Q_0S^0\), let
\(f_j:S^{d+j}\to QS^j\) be its successive adjoints and put
\(z_j=h(f_j)\). Suspension gives \(z_{j+1}=\sigma_*z_j\), and the
homology degree bounds imply that \(z_j=0\) for \(j>d\).
Thus a nonzero initial image has a last nonzero term \(z_e\).
We show that this terminal class is a square \(z_e=w^2\) whose
primitive, Steenrod-annihilated root still has nonzero suspension.
At a positive terminal level, the preceding spherical suspension forces
the root to have odd degree. At level zero, we instead exclude fourth powers
using spectrum pairings and integral characteristic classes
(Section~\ref{sec:powers}). The mapping cone of the next adjoint then
converts the square into a nonzero Steenrod-module extension
(Section~\ref{sec:cones}).

At positive level \(n\), weight counts the inputs of an extended power.
The primitives of weight \(2^k\), \(k\ge1\),
are dual to \(M_k(n)=\Sigma^n d_0^nD_k\), where \(D_k\) is a
binary Dickson algebra and \(d_0\) is its generator of degree
\(2^k-1\); the bottom summand is \(M_0(n)=\Sigma^n\F_2\).
Section~\ref{sec:weights} constructs these modules, identifies
the dual of homology suspension with inclusion of the next \(d_0\)-ideal,
and supplies compatible finite weight projections at level zero.
Section~\ref{sec:algebra} shows that the cone extension on a nonbottom
weight, when its detecting functional lifts from the preceding
suspension level, is possible only for \(n=k=1\) and in a
power-of-two degree. Cartier halving and an injection on even-degree
Tor cycles retain the exponent-floor information needed for this
restriction.

Section~\ref{sec:conclusion} assembles the argument. At a positive
terminal level the obstruction leaves only the bottom weight, giving
stable Hopf detection. At terminal level zero it forces the root to
have leading weight two, so the original spherical class has leading
weight four. Kuhn's theorem bounds its Adams filtration by two, and
Browder's identification gives Kervaire detection. Additivity of the
detecting invariants then gives the asserted span.

Figure~\ref{fig:proof-route} summarizes the two terminal cases.
\begin{figure}[t]
\centering
\small
\fbox{\begin{minipage}{0.88\linewidth}
\centering
An actual sphere map \(f_0:S^d\to Q_0S^0\), with \(h(f_0)\ne0\)\\[3pt]
\(z_j=\sigma_*^jh(f_0)\), with \(e=\max\{j:z_j\ne0\}\)
\end{minipage}}
\par\smallskip
\(\Downarrow\)\quad suspension kernel; the root survives by\\[2pt]
parity when \(e\ge1\), or exclusion of fourth powers when \(e=0\)
\par\smallskip
\fbox{\begin{minipage}{0.88\linewidth}
\centering
\(z_e=w^2\), with \(w\) primitive and Steenrod-annihilated,
\(\sigma_*w\ne0\)\\[3pt]
The cone of \(f_{e+1}\) supplies a nonzero square operation.
\end{minipage}}
\par\smallskip
\(\Downarrow\)\quad the algebra obstruction with the lifted functional
\par\smallskip
\begin{tabular}{@{}c@{\hspace{0.04\linewidth}}c@{}}
\fbox{\begin{minipage}{0.40\linewidth}\centering
\(e\ge1\)\\[3pt]
Only the bottom weight remains.\\
Stable Hopf detection\\
\(d=1,3,7\)
\end{minipage}}
&
\fbox{\begin{minipage}{0.40\linewidth}\centering
\(e=0\)\\[3pt]
The leading weight of \(z_0\) is four.\\
Adams filtration two\\
Kervaire detection
\end{minipage}}
\end{tabular}
\caption{The proof follows the successive adjoints of one stable
sphere class. The root \(w\) need not itself be spherical: the
next adjoint and its cone provide the realization information used
by the Steenrod-module obstruction.}
\label{fig:proof-route}
\end{figure}


\paragraph{Conventions.}
Spaces are based and of CW homotopy type; for a based space \(X\),
write \(QX=\Omega^\infty\Sigma^\infty X\).
Homology and cohomology have coefficients in \(\F_2\) unless stated
otherwise. Let \(A\) be the mod-2 Steenrod algebra and \(I=A^{>0}\)
its augmentation ideal. The dual homology action is denoted
\(\Sq_*^i\). A homology class is \(A\)-annihilated if every
\(\Sq_*^i\), \(i>0\), kills it. Equivalently, its functional on
cohomology vanishes on the image of \(I\). For a positive graded
cohomology algebra \(B\), write \(Q B=B/B^2\) for its
indecomposables.

Primitivity refers to the ordinary space diagonal, with the basepoint
as coaugmentation, unless the linear coproduct is explicitly specified.
Products in homology of an infinite loop
space are its additive Pontryagin products, written \(*\) when a
distinction is needed. The block pairing introduced later is written
\(\circ\). Homology suspension between infinite loop levels is
\(\sigma_*\). We distinguish it from stable homology evaluation
\(H_*(\Omega^\infty E)\to H_*(E)\), induced by the counit.

\begin{samepage}
The symbol \(S\) denotes the sphere spectrum. We write
\(D_rT=(T^{\wedge r})_{h\Sigma_r}\) for an extended power and
\(\mathbb P(T)=\bigvee_{r\ge0}D_rT\) for the free commutative ring
spectrum on \(T\).
All constructions are derived and homotopy coherent, including free
algebras, pushouts, group completion, localization and parametrized
families of maps. We use the usual mod-2 Dyer--Lashof operations on
infinite loop spaces and commutative ring spectra
\cite{CLM1976,BMMS1986}. Whenever a module is written with a suspension
shift, polynomial exponents refer to the unshifted monomial; its total
degree also includes that shift.
\par
\end{samepage}
